\section{Proofs of Section \ref{sec:HardFamilies}}\label{app:HardFamilies}

In this section, we restate and prove the theorems and propositions from Section \ref{sec:HardFamilies} in the order in which they appear.

\begin{proposition}\label{app:GraphClassGProps}
    For any  $\kay \geq 3$ and any $\ess \in \N$ family $\{(\GraphG_\ess^\kay,\, \GraphH_\ess^\kay)\}_{\ess \in \N}$, satisfies the following properties.
    \begin{enumerate}
        \item $|\VertexSet(\GraphG_\ess^\kay)|= |\VertexSet(\GraphH_\ess^\kay)|$ \label{GraphClassGPropsVertexNum}
        %$|\VertexSet(\GraphG_\ess^\kay)|= |\VertexSet(\GraphH_\ess^\kay)| = 2\kay \cdot (\ess + 1)$
        \item $\reg(\GraphG_\ess^\kay)=\reg(\GraphH_\ess^\kay)$\label{GraphClassGPropsRegularity}
        %$\reg(\GraphG_\ess^\kay)=\reg(\GraphG_\ess^\kay)=3\cdot(\kay +1)$, for $\ess \geq 2$
        \item $\omega(\GraphG_\ess^\kay) = \kay$ \label{GraphClassGPropsCliqueNum}
        \item $\chi(\GraphH_\ess^\kay) = \kay-1$ \label{GraphClassGPropsChromaticNum}
    \end{enumerate}
\end{proposition}
\begin{proof}
     \textbf{\ref{GraphClassGPropsVertexNum}.}
     Graph $\GraphG_0^\kay$ and its rewiring $\GraphH_0^\kay$ have the same number of vertices since no vertices are added or subtracted under rewiring. Any product of each of these graphs with a cycle $\Cycle_\enn$ preserves this equality.
    
    \textbf{\ref{GraphClassGPropsRegularity}.}
    Let graph vertex sets $\VertexSet_{K'}=\{\vertexv_0, \ldots, \vertexv_{\kay-1}\}$ and $\VertexSet_{K''}=\{\vertexv_\kay, \ldots, \vertexv_{2\kay-1}\}$ respectively make up the vertices of each disjoint $\CompleteGraph_\kay$ subgraph, $K'$ and $K''$, of $\GraphG_0^\kay$.
    Under rewiring of $\GraphG_0^\kay$, for exactly one edge $(\vertexv_i,\vertexv_j)$ of clique subgraph $K'$, its endpoint $\vertexv_j$ is changed to $\vertexv_{j+k}$ and for corresponding edge $(\vertexv_{i+k},\vertexv_{j+k})$ of clique subgraph $K''$ endpoint $\vertexv_{j+k}$ is changed to $\vertexv_{j}$. In this process $\vertexv_j$ loses and gains one neighbor, so regularity is preserved and equality $\reg(\GraphG_0^\kay)=\reg(\GraphH_0^\kay)$ holds. Any product of each of these graphs with a cycle $\Cycle_\enn$ preserves this equality by point \ref{GraphProdReg} of Proposition \ref{prop:SabidussiCombinedGeneral}.
    
    \textbf{\ref{GraphClassGPropsCliqueNum}.} 
    Since each $\GraphG_\ess^\kay$ is the product of $\CompleteGraph_\kay$ and two cycles, point \ref{GraphClassGPropsCliqueNum} follows from point \ref{ProdCliqInvsKC} of Proposition \ref{obs:ProdCliqInvs}.

    \textbf{\ref{GraphClassGPropsChromaticNum}.}
    We assume without loss of generality that edges $(\vertexv_0, \vertexv_1)$ and $(\vertexv_\kay, \vertexv_{\kay+1})$ of $\GraphG_0^\kay$ are rewired under $\GraphH_0^\kay$. Let coloring $\coloring: \VertexSet(\GraphH_0^\kay) \rightarrow \dbset{\kay-1}$ be defined such that $\coloring(\vertexv_0) = \coloring(\vertexv_1) = 0$ and $\coloring(\vertexv_i) = i-1$ for $i \in \{2, \ldots, \kay-1\}$. This coloring properly colors vertices in $\VertexSet_{K'}$ since vertices, $\vertexv_0$ and $\vertexv_1$ are non adjacent in $\GraphH_0^\kay$. We now extend coloring $\coloring$ to assign colors $\coloring(\vertexv_{\kay}) = \coloring(\vertexv_{\kay+1}) = \kay-2$ and $\coloring(\vertexv_{\kay + i}) = i-2$ for $i \in \{2, \ldots, \kay-1\}$. The colors of vertices $\vertexv_0,\vertexv_1$ do not conflict with their neighbors $\vertexv_{\kay+1},\vertexv_{\kay + i}$ in $\VertexSet_{K''}$ since $\coloring(\vertexv_0) = \coloring(\vertexv_1) = 0 \neq \kay-2 =\coloring(\vertexv_{\kay}) = \coloring(\vertexv_{\kay+1})$. More generally, for $i \in \{2, \ldots, \kay-1\}$ vertex $\vertexv_i$ does not conflict with its neighbor $\vertexv_{\kay+i}$ since $\coloring(\vertexv_i) = i-1\neq i-2 = \coloring(\vertexv_{\kay + i})$.

    By Proposition \ref{prop:SabidussiCombinedGeneral} it holds that $\chi(\GraphH_\ess^\kay) = \kay-1$ for all $\ess \in \N$.
\end{proof}

\begin{theorem}
    Let $\mathcal{N}$ be a message-passing graph neural network whose aggregation and update
functions are injective. Then for any pair of graphs $(\GraphG, \GraphH)$ of family  $\{(\GraphG_\ess^\kay,\, \GraphH_\ess^\kay)\}_{\ess \in \N}$, $\GNN(\GraphG) = \GNN(\GraphH)$.
\end{theorem}
\begin{proof}
    By Proposition \ref{prop:GraphClassGProps} graphs $\GraphG$ and $\GraphH$ are on the same number of vertices and have the same regularity. Hence, $\GNN(\GraphG) = \GNN(\GraphH)$ by Observation \ref{obs:WL-regular-undistinguished} and Proposition \ref{prop:GNN-iff-WL}.
\end{proof}


\begin{theorem}
    Let $\mathcal{N}_2$ be a 2-GNN whose aggregation and update
    functions are injective. Then for any pair $(\GraphG,\GraphH)$ in the family
    $\{(\GraphG_\ess^\kay,\, \GraphH_\ess^\kay)\}_{\ess \in \N}$, we have $\mathcal{N}_2(\GraphG)=\mathcal{N}_2(\GraphH)$.
\end{theorem}
% \begin{proof}
%     We show that 2-WL stabilizes in two steps. That is, we show that for any pair $(\GraphG,\GraphH)$ in $\{(\GraphG_\ess^\kay,\, \GraphH_\ess^\kay)\}_{\ess \in \N}$, the cardinalities of initial coloring $\coloring^{(0)}$ and consecutive coloring $\coloring^{(1)}$ are equal. 

%     Without loss of generality we assume that initial coloring $\coloring^{(0)}$ assigns to each edge color $\coloring_1$ and each non-edge color $\coloring_2$ where $\coloring_1 \neq \coloring_2$. By points \ref{GraphClassGPropsVertexNum} and \ref{GraphClassGPropsRegularity} of \ref{prop:GraphClassGProps} graphs $\GraphG$ and $\GraphH$ have the same number of edges and non-edges. Let $\edgee$ and $\nonedgee$ respectively be an edge and non-edge of $\EdgeSet(\GraphG) \cup \EdgeSet(\GraphH)$. We show that coloring $\coloring^{(1)}$ assigns color $\coloring_3$ to each edge $\edgee$ and $\coloring_4$ to each non-edge $\nonedgee$, where $\coloring_3 \neq \coloring_4$.

%     Let edge $\edgee = (\vertexu, \vertexv)$ be chosen arbitrarily from $\EdgeSet(\GraphG) \cup \EdgeSet(\GraphH)$. The set of neighbors of $\edgee$ is exactly the union of edges of $\GraphG$ ($\GraphH$) and non-edges of of $\GraphG$ ($\GraphH$) with one endpoint in $\edgee$ and one endpoint outside of $\edgee$. The number of edges incident to endpoint $\vertexu$ ($\vertexv$) is therefore given by $\reg-1$ (we do not count $\edgee$ itself) where $\reg$ is the regularity of $\GraphG$ ($\GraphH$)(recall that $\GraphG$ and $\GraphH$ have the same regularity of point \ref{GraphClassGPropsRegularity} of \ref{prop:GraphClassGProps}). Meanwhile, the number of non-edges incident to $\vertexu$ ($\vertexv$) is given by $|\VertexSet(\GraphG)| - \reg$ ($|\VertexSet(\GraphH)| - \reg$) (recall that vertex sets $\VertexSet(\GraphG)$ and $\VertexSet(\GraphH)$ have the same cardinality by \ref{GraphClassGPropsVertexNum} of \ref{prop:GraphClassGProps}). Hence and arbitrarily chosen $\edgee$ of $\EdgeSet(\GraphG) \cup \EdgeSet(\GraphH)$ has the same neighborhood and is assigned the same color under $\coloring^{(1)}$.

%     An analogous argument holds for arbitrarily chosen non-edges $\nonedgee$ of $\EdgeSet(\GraphG) \cup \EdgeSet(\GraphH)$.

%     Since color assignments under $\coloring^{1}$ are completely determined by isomorphism type, the image of $\coloring^{1}$ has the same cardinality as the image of $\coloring^{0}$.

    
% \end{proof}
\begin{proof}
We show that the $2$-WL refinement stabilizes after one refinement step. That is, for any pair
$(\GraphG,\GraphH)$ in the family
$\{(\GraphG_\ess^\kay,\, \GraphH_\ess^\kay)\}_{\ess \in \N}$,
the images of the initial coloring $\coloring^{(0)}$ and the refined coloring
$\coloring^{(1)}$ have the same cardinality.

Without loss of generality, assume that the initial coloring $\coloring^{(0)}$ assigns a color
$\coloring_1$ to each edge and a distinct color $\coloring_2 \neq \coloring_1$ to each non-edge.
By points~\ref{GraphClassGPropsVertexNum} and~\ref{GraphClassGPropsRegularity} of
Proposition~\ref{prop:GraphClassGProps}, the graphs $\GraphG$ and $\GraphH$ have the same number
of vertices, the same regularity, and consequently the same number of edges and non-edges.

Let $\edgee = (\vertexu,\vertexv)$ be an arbitrary edge in
$\EdgeSet(\GraphG) \cup \EdgeSet(\GraphH)$. In the $2$-WL refinement, the neighbors of $\edgee$
are precisely the $2$-sets that share exactly one vertex with $\edgee$. These consist of all edges
and non-edges incident to either $\vertexu$ or $\vertexv$, excluding $\edgee$ itself.

Since both $\GraphG$ and $\GraphH$ are $\reg$-regular, each endpoint $\vertexu$ and $\vertexv$
is incident to exactly $\reg - 1$ further edges. Moreover, as both graphs have the same number of
vertices, the number of non-edges incident to each endpoint is
$|\VertexSet(\GraphG)| - \reg - 1$ (and likewise for $\GraphH$).
Thus, every edge in $\EdgeSet(\GraphG) \cup \EdgeSet(\GraphH)$ has an identical multiset of
neighbor colors under $\coloring^{(0)}$, and is therefore assigned the same color, say
$\coloring_3$, under $\coloring^{(1)}$.

An analogous argument applies to any non-edge
$\nonedgee \notin \EdgeSet(\GraphG) \cup \EdgeSet(\GraphH)$, showing that all non-edges receive a
common color $\coloring_4 \neq \coloring_3$ under $\coloring^{(1)}$.

It follows that the refined coloring $\coloring^{(1)}$ distinguishes edges from non-edges but
introduces no further refinement. Hence, the image of $\coloring^{(1)}$ has the same cardinality
as the image of $\coloring^{(0)}$, and the $2$-WL refinement stabilizes after one step.

By Proposition~\ref{prop:kGNN-iff-kWL}, $2$-GNNs with injective aggregation and update functions
cannot distinguish $\GraphG$ from $\GraphH$, and therefore
$\mathcal{N}_2(\GraphG) = \mathcal{N}_2(\GraphH)$.
\end{proof}


\begin{proposition}\label{app:Dis3GNN}
Let $\mathcal{N}_3$ be a 3-GNN. There exist aggregation and update functions such that for any two graphs $\{\GraphG_0^\kay\}$ and $\{\GraphH_0^\kay\}$,
\[
\mathcal{N}(\GraphG_0^\kay) \neq \mathcal{N}(\GraphH_0^\kay),
\]    
for any $k \geq 3$.
\end{proposition}
\begin{proof}
    See that for each $k\geq3$ at least one induced triangle of $\GraphG_0^\kay$ is removed under rewiring to $\GraphH_0^\kay$. Graphs $\{\GraphG_0^\kay\}$ and $\{\GraphH_0^\kay\}$ are therefore distinguished in a single 3-WL refinement step since these graphs will have differing color-wise distribution of $3$-sets. Next see that taking graph products with even cycles does not introduce new triangles and therefore our argument holds for every family $\Fam_\kay$. By Proposition \ref{prop:kGNN-iff-kWL} there must exist $\mathcal{N}_3$ distinguishing each graph family $\Fam_\kay$.
\end{proof}
